% =============================================================================
% CHAPTER 1 — INTRODUCTION GÉNÉRALE
% =============================================================================

\chapteropening{01}{Introduction\\Générale}{Introduction Générale}{chap:introduction}

Ce rapport présente le travail réalisé dans le cadre du projet de fin d'études effectué au sein de Djagora. Le stage a porté sur l'accompagnement entrepreneurial de startups, l'étude de modèles économiques, et le développement de produits minimaux viables (MVP).

\vspace{12pt}

Le présent rapport est structuré en trois chapitres :

\begin{itemize}[leftmargin=*]
  \item \textbf{Chapitre 1} présente le cadre général du stage, l'organisme d'accueil et l'écosystème entrepreneurial.
  \item \textbf{Chapitre 2} expose le cadre théorique : startup, Business Model Canvas, Lean Canvas, métriques et méthodologie SCRUM.
  \item \textbf{Chapitre 3} détaille la réalisation pratique à travers les études de cas des startups accompagnées.
\end{itemize}

\vspace{12pt}

Ce projet de fin d'études s'inscrit dans une démarche globale visant à concevoir une plateforme intelligente de pilotage de l'employabilité des étudiants. À travers ce rapport, nous présentons non seulement les fondamentaux du contexte académique et entrepreneurial, mais aussi l'approche systémique qui sous-tend la solution LogiLink. Notre objectif est de démontrer comment une approche méthodique, associée à des technologies avancées d'intelligence artificielle, peut contribuer à réduire le décalage entre la formation académique et les exigences réelles du marché de l'emploi, particulièrement dans les secteurs du transport et de la logistique.

\vspace{12pt}

Les chapitres qui suivent développent cette vision en détail, en s'appuyant sur une analyse rigoureuse des besoins, une spécification précise des exigences, et une implémentation guidée par la méthodologie Agile SCRUM.
